\section{Conclusion}
\label{sec:conclusion}

We presented VGuide, a CE-guided LLM predicate oracle for CEGAR-based software
verification. VGuide intercepts the CEGAR refinement loop at the first spurious
counterexample, uses \LLMName{} (non-thinking, \LLMLatencyMedian{}\,s median latency)
to propose candidate predicates from a structured context pack, and injects only
predicates that pass a three-tier validation cascade (contract, vocabulary, SMT
entailment). The system maintains the soundness invariant of predicate abstraction
unconditionally: no LLM output reaches the abstract domain without formal certification.
On 764 SV-COMP Loops ReachSafety tasks, VGuide gains \LoopsDelta{} tasks over
base predicate analysis and \PortfolioLoopsDelta{} tasks over the svcomp26 portfolio;
combined competition-grade evaluation shows \CompTotalDelta{} additional solved tasks
with zero wrong verdicts. An ablation confirms that CE context is necessary---source-code-only
LLM queries yield no improvement whatsoever.

Future work includes:
\begin{itemize}
  \item \textbf{FALSE-path guidance.} Extending VGuide with a predicate injection
        mode oriented toward counterexample witness discovery, improving the
        \textsc{False} verdict count on bug-finding tasks.
  \item \textbf{Offline and local LLM deployment.} Pre-computing predicate suggestions
        before the competition isolation window, or fine-tuning a local model on
        curated CEGAR-context data, to enable VGuide deployment in network-isolated
        environments such as SV-COMP.
  \item \textbf{Structured loop invariant templates.} Using program-structure analysis
        to generate predicate templates (e.g., affine constraints between specific
        variable pairs) that guide the LLM toward more precise and less noisy
        suggestions, reducing the failure mode observed in \texttt{nested-3}.
\end{itemize}
